\chapter{Rangkaian Listrik dan Keselamatan}\label{ch:g6:circuits-and-safety}

Tiga tahun bergelut dengan lingkar, \omterm{def:g3:first-electric-circuit:bulb}{bohlam}, dan \omterm{def:g5:series-parallel-circuits:parallel}{cabang} sudah
menjadikanmu pembangun rangkaian. Tahun ini kamu menjadi sesuatu yang
lebih langka: \emph{penulis} rangkaian. Tukang listrik di setiap
negeri membaca dan menggambar bahasa tanpa kata yang sama --- dan tata
bahasanya, ditambah satu bab kearifan keselamatan yang mahal harganya,
itulah yang berdiri di antara pemasangan yang cerdik dan jari yang
melepuh.

\section{Bahasa tanpa kata}

\begin{definition}[Diagram rangkaian]\label{def:g6:circuits-and-safety:diagram}
Sebuah \emph{diagram rangkaian}\index{diagram rangkaian} adalah gambar
baku sebuah rangkaian: tiap alat ditampilkan oleh
\emph{lambang}\index{lambang} yang disepakati untuknya, dan kawatnya
oleh garis lurus yang digambar bersudut siku. Diagram menunjukkan
bagaimana rangkaiannya \emph{terhubung} --- siapa memberi makan siapa,
apa bercabang di mana --- dan dengan sengaja mengabaikan bagaimana
kawat yang sungguhan meliuk di atas meja.
\end{definition}

\begin{omfigure}
\begin{tikzpicture}[scale=0.8]
  \draw (0.4,3.4) to[battery1] (2.4,3.4);
  \node[right, font=\small] at (2.7,3.4)
    {baterai (garis panjang: $+$)};
  \draw (0.4,2.2) to[lamp] (2.4,2.2);
  \node[right, font=\small] at (2.7,2.2) {lampu};
  \draw (0.4,1.0) to[spst] (2.4,1.0);
  \node[right, font=\small] at (2.7,1.0) {sakelar, terbuka};
  \draw (7.4,3.4) to[cspst] (9.4,3.4);
  \node[right, font=\small] at (9.7,3.4) {sakelar, tertutup};
  \draw (7.4,2.2) -- (8.05,2.2);
  \draw[thick] (8.4,2.2) circle (0.35);
  \node[font=\small] at (8.4,2.2) {M};
  \draw (8.75,2.2) -- (9.4,2.2);
  \node[right, font=\small] at (9.7,2.2) {motor};
  \draw (7.4,1.0) -- (8.05,1.0);
  \draw[thick] (8.4,1.0) circle (0.35);
  \draw[thick] (8.16,0.76) -- (8.64,1.24);
  \draw (8.75,1.0) -- (9.4,1.0);
  \node[right, font=\small] at (9.7,1.0) {bel listrik};
\end{tikzpicture}

{\small Abjadnya: enam \omterm{def:g6:circuits-and-safety:diagram}{lambang} membawa hampir setiap rangkaian dalam
buku ini. Tanda yang sama dibaca dari bengkel ke bengkel di seluruh
dunia.}
\end{omfigure}

\begin{method}[Dari meja ke kertas, dan sebaliknya]\label{met:g6:circuits-and-safety:translate}
Untuk \emph{menggambar} rangkaian yang sudah dibangun:
\begin{enumerate}
  \item daftarkan alatnya, lalu pilih \omterm{def:g6:circuits-and-safety:diagram}{lambang} masing-masingnya;
  \item ikuti lingkar yang sungguhan dengan jari mulai dari kutub $+$
        \omterm{def:g3:first-electric-circuit:battery}{baterai}, sambil mencatat urutan alatnya dan setiap
        percabangannya;
  \item gambarlah ulang sebagai persegi panjang yang rapi: \omterm{def:g6:circuits-and-safety:diagram}{lambangnya}
        berjarak, kawatnya lurus, sudutnya siku --- hubungan yang
        sama, tanpa liukan.
\end{enumerate}
Untuk \emph{membangun} dari sebuah diagram, jalankan penerjemahannya
mundur --- lalu sesudahnya, periksalah seperti biasa: setiap lingkar
yang ditelusuri dari $+$ harus mencapai $-$ lewat setidaknya satu alat.
\end{method}

\begin{example}[Rangkaian yang sama, dua potret]\label{ex:g6:circuits-and-safety:portraits}
Di atas meja: sebuah \omterm{def:g3:first-electric-circuit:battery}{baterai}, sekusut kawat berjepit buaya yang
melingkar di atas kotak pensil, sebuah lampu yang bersandar pada buku,
sebuah \omterm{ex:g3:first-electric-circuit:switch}{sakelar} yang menjuntai. Di atas kertas: sebuah persegi panjang
yang tenang --- \omterm{def:g3:first-electric-circuit:battery}{baterai} di kiri, \omterm{ex:g3:first-electric-circuit:switch}{sakelar} di atas, lampu di kanan.
Kekusutan dan persegi panjang itu adalah \emph{rangkaian yang sama}:
setiap hubungannya sama persis. Untuk memutuskan apakah dua rangkaian
itu sama, ajukan satu pertanyaan saja: siapa terhubung dengan siapa?
\end{example}

\begin{example}[Motor dan bel ikut bermain]\label{ex:g6:circuits-and-safety:motor}
Dua pendatang baru dalam kotak alatmu berperilaku persis seperti lampu
di dalam diagram: \emph{motor}\index{motor} (\omterm{def:g6:circuits-and-safety:diagram}{lambang}: huruf M di dalam
lingkaran) berputar setiap kali arus mengalir melewatinya --- kipas,
mobil mainan, bor; \emph{bel listrik}\index{bel listrik} (lingkaran
yang terpalang) berbunyi setiap kali diberi makan --- bel pintu,
alarm. Aturan seri dan paralel berlaku bagi keduanya tanpa perubahan:
\omterm{ex:g3:first-electric-circuit:switch}{sakelar} yang seri memerintah mereka; \omterm{def:g5:series-parallel-circuits:parallel}{cabang} yang paralel membuat
mereka mandiri.
\end{example}

\section{Aturan lama dalam bahasa yang baru}

\begin{example}[Membaca diagram seperti ahli]\label{ex:g6:circuits-and-safety:reading}
Terhadap diagram apa pun, tanyailah ia dengan jalan-jari dari bab
\omterm{def:g5:series-parallel-circuits:parallel}{rangkaian paralel}: dari $+$, jalanilah setiap jalan menuju $-$. Semua
alat pada satu jalan: seri --- satu celah menghentikan semuanya,
urutannya bebas, kecerahannya dibagi. Jalan yang bercabang pada titik
sambungan: paralel --- \omterm{def:g5:series-parallel-circuits:parallel}{cabangnya} mandiri, masing-masing terang penuh,
dan penempatan \omterm{ex:g3:first-electric-circuit:switch}{sakelarnya} kini menjadi keputusan yang sungguhan.
Bahasanya berubah; hukumnya tidak.
\end{example}

\begin{omfigure}
\begin{tikzpicture}[scale=0.8]
  \draw (0,0) to[battery1] (0,3.2)
    to[cspst] (2.6,3.2);
  \draw (2.6,3.2) to[short] (2.6,2.4)
    to[lamp] (5.8,2.4) to[short] (5.8,3.2);
  \draw (2.6,3.2) to[short] (2.6,4.0) -- (3.4,4.0);
  \draw[thick] (3.75,4.0) circle (0.35);
  \node[font=\small] at (3.75,4.0) {M};
  \draw (4.1,4.0) -- (5.8,4.0) to[short] (5.8,3.2);
  \draw (5.8,3.2) to[short] (7.4,3.2) to[short] (7.4,0)
    to[short] (0,0);
  \fill (2.6,3.2) circle (2pt);
  \fill (5.8,3.2) circle (2pt);
\end{tikzpicture}

{\small Bacalah sebelum membangunnya: satu \omterm{ex:g3:first-electric-circuit:switch}{sakelar} pada jalan
bersamanya, lalu sebuah percabangan --- lampu pada satu \omterm{def:g5:series-parallel-circuits:parallel}{cabang}, \omterm{ex:g6:circuits-and-safety:motor}{motor}
kipas pada \omterm{def:g5:series-parallel-circuits:parallel}{cabang} yang lain. \omterm{ex:g3:first-electric-circuit:switch}{Sakelarnya} memerintah keduanya;
\omterm{def:g5:series-parallel-circuits:parallel}{cabang-cabangnya} mengabaikan kesulitan satu sama lain.}
\end{omfigure}

\section{Bab tentang bekas luka}

Aturan keselamatan listrik dibayar dengan kebakaran dan hal yang lebih
buruk. Inilah aturannya, beserta alasannya --- seorang profesional
mengetahui keduanya.

\begin{definition}[Korsleting]\label{def:g6:circuits-and-safety:short}
Yang disebut \emph{korsleting}\index{korsleting} terjadi ketika sebuah
jalan penghantar yang telanjang menghubungkan kedua sisi \omterm{def:g3:first-electric-circuit:battery}{baterai} (atau
stopkontak) secara \emph{langsung}, dengan melewatkan setiap alat.
Karena tidak ada yang berguna menghalangi jalannya, arusnya menyerbu
liar --- kawat memanas, isolasinya melebur, \omterm{def:g3:first-electric-circuit:battery}{baterai} rusak, dan
pemasangan listrik rumah dapat terbakar. Bab rangkaian tahun depan
mempelajari penjahat ini dari dekat; tahun ini, kenalilah wajahnya dan
jauhkan ia dari bangunanmu.
\end{definition}

\begin{remark}[Para pengawal rumah]\label{rem:g6:circuits-and-safety:guards}
Lihatlah ke dalam kotak sekring di rumah: berderet \omterm{ex:g3:first-electric-circuit:switch}{sakelar} kecil ---
\emph{pemutus arus} --- satu untuk tiap \omterm{def:g5:series-parallel-circuits:parallel}{cabang} \omterm{def:g5:series-parallel-circuits:parallel}{rangkaian paralel} besar
milik rumahmu. Masing-masing berjaga atas \omterm{def:g5:series-parallel-circuits:parallel}{cabangnya}: kalau sebuah
kerusakan membiarkan arusnya menyerbu --- \omterm{def:g6:circuits-and-safety:short}{korsleting}, stopkontak yang
kelebihan beban --- sang pengawal mematikan dirinya sendiri dalam
sekejap, memutus \omterm{def:g5:series-parallel-circuits:parallel}{cabangnya} sebelum kawatnya terlalu panas. Ketika
sebuah pemutus arus ``turun'', ia bukan sedang berulah: ia baru saja
melakukan satu-satunya tugasnya. Carilah kerusakannya sebelum
menaikkannya kembali.
\end{remark}

\begin{remark}[Aturannya, beserta alasannya]\label{rem:g6:circuits-and-safety:rules}
\begin{enumerate}
  \item \emph{Percobaan dengan \omterm{def:g3:first-electric-circuit:battery}{baterai} saja; stopkontak, jangan
        pernah.} Listrik rumah mendorong ratusan kali lebih kuat ---
        lewat kulit yang lembap ia dapat menghentikan jantung.
  \item \emph{Jangan ada kawat telanjang yang melintang langsung
        pada \omterm{def:g3:first-electric-circuit:battery}{baterai}.} Itulah \omterm{def:g6:circuits-and-safety:short}{korsleting}: sel yang rusak, kawat yang
        panas, jari yang melepuh.
  \item \emph{Air dan listrik rumah tidak pernah berbagi jangkauan
        satu ruangan:} tidak ada peranti di dekat bak mandi, tidak
        ada tangan basah pada steker --- kulit yang lembap adalah
        penghantar.
  \item \emph{Jangan pernah menyodokkan apa pun ke dalam
        stopkontak.} Pengawal di kotak sekring memang cepat, tetapi
        kamu tidak boleh pernah bertaruh padanya.
  \item \emph{Kabel yang rusak masuk ke tempat sampah,} bukan ke
        laci: isolasi yang retak adalah pagar yang berlubang.
\end{enumerate}
\end{remark}

\section{Latihan}

\begin{exercise}[$\star$]\label{exo:g6:circuits-and-safety:1}
Apa yang ditunjukkan sebuah \omterm{def:g6:circuits-and-safety:diagram}{diagram rangkaian} dengan setia, dan apa
yang diabaikannya dengan sengaja?
\end{exercise}

\begin{exercise}[$\star$]\label{exo:g6:circuits-and-safety:2}
Gambarlah dari ingatan \omterm{def:g6:circuits-and-safety:diagram}{lambang} berikut: \omterm{def:g3:first-electric-circuit:battery}{baterai}; \ lampu; \ \omterm{ex:g3:first-electric-circuit:switch}{sakelar}
terbuka; \ \omterm{ex:g3:first-electric-circuit:switch}{sakelar} tertutup; \ \omterm{ex:g6:circuits-and-safety:motor}{motor}; \ \omterm{ex:g6:circuits-and-safety:motor}{bel listrik}. Tanda mana yang
menyebutkan sisi $+$ \omterm{def:g3:first-electric-circuit:battery}{baterai}?
\end{exercise}

\begin{exercise}[$\star$]\label{exo:g6:circuits-and-safety:3}
Gambarlah diagram sebuah bel pintu: \omterm{def:g3:first-electric-circuit:battery}{baterai}, \omterm{ex:g3:first-electric-circuit:switch}{sakelar} tekan, dan \omterm{ex:g6:circuits-and-safety:motor}{bel
listrik} pada satu lingkar. Mengapa belnya berhenti ketika jarinya
diangkat?
\end{exercise}

\begin{exercise}[$\star$]\label{exo:g6:circuits-and-safety:4}
Dua rangkaian dibangun dari kotak suku cadang yang sama: apakah
keduanya ``rangkaian yang sama''? Satu pertanyaan apa yang memutuskan?
\end{exercise}

\begin{exercise}[$\star$]\label{exo:g6:circuits-and-safety:5}
Pada diagram lampu-dan-motor dalam bab ini, apa yang terjadi pada
kipasnya kalau lampunya putus? Dan pada keduanya kalau \omterm{ex:g3:first-electric-circuit:switch}{sakelarnya}
terbuka? Aturan lama yang mana yang menjawab?
\end{exercise}

\begin{exercise}[$\star$]\label{exo:g6:circuits-and-safety:6}
Apa itu \omterm{def:g6:circuits-and-safety:short}{korsleting}, dan mengapa ia memanaskan kawat? Sebutkan dua
korbannya.
\end{exercise}

\begin{exercise}[$\star$]\label{exo:g6:circuits-and-safety:7}
Sebuah pemutus arus di kotak sekring baru saja turun. Apa yang baru
saja dilakukannya untuk rumah itu? Apa yang harus dilakukan sebelum
menaikkannya kembali?
\end{exercise}

\begin{exercise}[$\star$]\label{exo:g6:circuits-and-safety:8}
Sebutkan alasan di balik tiap aturan: tidak ada peranti di dekat bak
mandi; \ kabel yang rusak masuk tempat sampah; \ permainan \omterm{def:g3:first-electric-circuit:battery}{baterai}
boleh, stopkontak jangan pernah.
\end{exercise}

\begin{exercise}[$\star\star$]\label{exo:g6:circuits-and-safety:9}
Terjemahkan menjadi sebuah diagram (dengan kata atau gambar): sebuah
\omterm{def:g3:first-electric-circuit:battery}{baterai}; dari kutub $+$-nya, sebuah kawat menuju \omterm{ex:g3:first-electric-circuit:switch}{sakelar} tertutup;
dari \omterm{ex:g3:first-electric-circuit:switch}{sakelarnya} sebuah percabangan --- satu \omterm{def:g5:series-parallel-circuits:parallel}{cabang} lewat \omterm{ex:g6:circuits-and-safety:motor}{bel listrik},
satu lewat lampu; kedua \omterm{def:g5:series-parallel-circuits:parallel}{cabangnya} bergabung kembali lalu pulang ke
$-$. Ramalkan apa yang terdengar dan terlihat, dan apa yang berubah
kalau \omterm{def:g5:series-parallel-circuits:parallel}{cabang} belnya saja mendapat \omterm{ex:g3:first-electric-circuit:switch}{sakelar} tambahannya sendiri, dalam
keadaan terbuka.
\end{exercise}

\begin{exercise}[$\star\star$]\label{exo:g6:circuits-and-safety:10}
Sebuah kipas mainan dipasang: kutub $+$ \omterm{def:g3:first-electric-circuit:battery}{baterai} ke \omterm{ex:g6:circuits-and-safety:motor}{motor}, \omterm{ex:g6:circuits-and-safety:motor}{motor} ke
kutub $-$ \omterm{def:g3:first-electric-circuit:battery}{baterai}, dan --- ``supaya kukuh'' --- satu kawat telanjang
tambahan langsung dari $+$ ke $-$. Kipasnya nyaris tidak berputar dan
kawat tambahannya menjadi panas. Diagnosislah dengan
\cref{def:g6:circuits-and-safety:short}: jalan mana yang lebih disukai
arusnya, dan mengapa bangunannya berbahaya?
\end{exercise}

\begin{exercise}[$\star\star$]\label{exo:g6:circuits-and-safety:11}
Kebiasaan bengkel Kakek: sebelum menyentuh lampu apa pun, ia mematikan
\omterm{ex:g3:first-electric-circuit:switch}{sakelar} dindingnya \emph{dan} pemutus arus seluruh \omterm{def:g5:series-parallel-circuits:parallel}{cabangnya} di kotak
sekring. Dengan cara berpikir seri, jelaskan apa yang dicapai tiap
pemutusan itu --- dan mengapa orang yang cermat melakukan keduanya.
\end{exercise}

\begin{exercise}[$\star\star\star$]\label{exo:g6:circuits-and-safety:12}
Rancanglah pemasangan kabel sebuah panggung boneka: satu \omterm{ex:g3:first-electric-circuit:switch}{sakelar} utama
yang menggelapkan segalanya; dua lampu panggung yang harus dapat mati
sendiri-sendiri; sebuah \omterm{ex:g6:circuits-and-safety:motor}{motor} tirai dan sebuah bel tanda pertunjukan,
masing-masing dengan \omterm{ex:g3:first-electric-circuit:switch}{sakelar} pengendalinya sendiri. Ceritakan (atau
gambarlah) diagramnya --- alat mana yang menunggangi jalan bersama,
mana yang menunggangi \omterm{def:g5:series-parallel-circuits:parallel}{cabang}, di mana tiap \omterm{ex:g3:first-electric-circuit:switch}{sakelarnya} duduk --- lalu
periksalah rancanganmu terhadap setiap aturan seri dan paralel yang
kamu miliki.
\end{exercise}

\section{Soal: Memasang Kabel Gudang Kebun}

\begin{problem}[{Soal akhir pekan --- sebuah keluarga memasang kabel
gudang kebunnya; diagram sebelum obeng; daftar periksa sang
pemeriksa}]\label{pb:g6:circuits-and-safety:1}
Keluarga itu merencanakan jaringan kecil bertenaga \omterm{def:g3:first-electric-circuit:battery}{baterai} untuk
gudangnya di atas kertas lebih dahulu --- seperti yang dilakukan para
profesional.

\medskip\noindent\textbf{Bagian I --- Rencana di atas kertas.}
Daftar keinginannya: satu lampu langit-langit yang disakelari di
pintu; satu lampu meja kerja dengan \omterm{ex:g3:first-electric-circuit:switch}{sakelarnya} sendiri; satu kipas
pengudaraan kecil (\omterm{ex:g6:circuits-and-safety:motor}{motor}) dengan \omterm{ex:g3:first-electric-circuit:switch}{sakelarnya} sendiri.
\begin{enumerate}
  \item Mengapa ketiga alatnya harus menunggangi \omterm{def:g5:series-parallel-circuits:parallel}{cabang} paralel
        alih-alih satu lingkar seri? Sebutkan dua kenyamanan yang
        akan rusak kalau dipasang seri.
  \item Tiap \omterm{def:g5:series-parallel-circuits:parallel}{cabang} membawa \omterm{ex:g3:first-electric-circuit:switch}{sakelarnya} sendiri. Apa yang diperintah
        tiap \omterm{ex:g3:first-electric-circuit:switch}{sakelar}, dan apa yang tidak diperintah satu pun di
        antaranya?
  \item Di mana \omterm{ex:g3:first-electric-circuit:switch}{sakelar} utama harus duduk supaya menggelapkan dan
        membungkam seluruh gudangnya sekaligus?
  \item Gambarlah (atau ceritakan dengan tepat) diagram lengkapnya:
        \omterm{def:g3:first-electric-circuit:battery}{baterai}, \omterm{ex:g3:first-electric-circuit:switch}{sakelar} utama, tiga \omterm{def:g5:series-parallel-circuits:parallel}{cabang}, empat \omterm{ex:g3:first-electric-circuit:switch}{sakelar}
        seluruhnya.
\end{enumerate}

\medskip\noindent\textbf{Bagian II --- Bangunannya, diperiksa.}
\begin{enumerate}[resume]
  \item Sesudah dibangun dan dinyalakan, lampu langit-langitnya
        menyala --- tetapi redup, dan lampu meja kerjanya hanya
        menyala ketika \omterm{ex:g3:first-electric-circuit:switch}{sakelar} \emph{kipasnya} tertutup. Jalan-jari
        menemukan lampu meja dan kipasnya berbagi satu \omterm{def:g5:series-parallel-circuits:parallel}{cabang}, yang
        satu di belakang yang lain. Sebutkan susunannya, lalu
        jelaskan kedua gejalanya dengan aturan lama.
  \item Berikan perbaikan satu kawat yang memulihkan rencananya.
  \item Ketika berbenah, sebuah obeng sejenak terbaring melintang
        pada kedua kutub telanjang \omterm{def:g3:first-electric-circuit:battery}{baterainya} lalu menjadi hangat.
        Sebutkan nama peristiwanya, lalu nyatakan apa yang harus
        diperiksa keluarganya tentang \omterm{def:g3:first-electric-circuit:battery}{baterainya} sesudah itu.
  \item Orang tua yang bertugas memeriksa menambahkan sebuah aturan:
        ``setiap sambungan dibungkus, tidak ada tembaga telanjang di
        mana pun.'' Dua bahaya yang mana yang dijaga oleh satu
        aturan itu?
\end{enumerate}

\medskip\noindent\textbf{Bagian III --- Malam badai.}
\begin{enumerate}[resume]
  \item Malam yang hujan: atap gudangnya bocor ke atas meja kerja.
        Mengapa keluarganya tidak boleh menyentuh pemasangan
        kabelnya dengan tangan basah, bahkan di gudang bertenaga
        \omterm{def:g3:first-electric-circuit:battery}{baterai} --- dan mengapa sentuhan yang sama akan jauh lebih
        gawat di rumah yang memakai listrik rumah?
  \item Lampu rumahnya sendiri berkedip dan satu pemutus arus turun.
        Apa yang dideteksi pemutus arus itu, dan bagaimana urutan
        tindakan yang bijak?
  \item Pagi hari: keluarganya berdebat tentang menambah bel pintu
        --- \omterm{ex:g6:circuits-and-safety:motor}{bel listrik} di gudang, \omterm{ex:g3:first-electric-circuit:switch}{sakelar} tekan di rumah, sepasang
        kawat panjang di antaranya. Seri atau paralel terhadap
        jaringan gudangnya? Dari mana arusnya datang? Sketsalah
        lingkarnya.
  \item Tulislah piagam keselamatan tiga baris milik keluarga itu
        untuk gudangnya, satu baris masing-masing: stopkontak, air,
        kawat telanjang.
\end{enumerate}
\end{problem}
